% !TEX root = ../main.tex

\section{一维教学模型}\label{sec:1d-model}

\subsection{模型定义}\label{subsec:model-def}

我们考虑一个一维晶格，上面 populate 了两种粒子：
\begin{itemize}
    \item $(+)$ 粒子：以恒定速率 $\lambda$ 向右移动到 $i+1$
    \item $(-)$ 粒子：以恒定速率 $\lambda$ 向左移动到 $i-1$
\end{itemize}

设晶格点 $i$ 上物种 $\pm$ 的粒子数为 $n_i^\pm$。系统的动力学由两个过程组成：

\begin{keybox}[模型动力学]
\begin{enumerate}
    \item \textbf{流动(Streaming)}：粒子以速率 $\lambda$ 从 $i$ 移动到 $i\pm 1$
    \item \textbf{转换(Conversion)}：粒子以速率 $r_\pm(n_i^+, n_i^-)$ 从 $+$ 转换为 $-$（或反之）
\end{enumerate}
\end{keybox}

转换率的选择反映了反排列相互作用的物理。在本文中，选择最简单的非线性形式（忽略自相互作用）：
\begin{equation}\label{eq:conversion-rate}
\boxed{r_\pm(n^+, n^-) = r_0 + (n^\pm - 1)r_1}
\end{equation}

其中 $r_0$ 是``自发转换率''，$r_1$ 是与局部磁化强度相关的反排列相互作用强度。

\begin{remark}[关于转换率形式的注释]
这里要求 $r_1$ 远大于 $r_0$，因为我们关注的是由于局部磁化引起的转换。转换率中只保留到 $O(n^+, n^-)$ 的项是关键的——更高阶的项会导致 Fokker-Planck 方程中出现高于二阶的导数，从而无法通过对应的 Langevin 方程直接求解。
\end{remark}

\subsection{Master 方程}\label{subsec:master-eq}

设 $P(\{n_i^+, n_i^-\})$ 为找到特定占据数集合的概率分布。Master 方程为：
\begin{equation}\label{eq:master-eq}
\partial_t P = -S^\lambda_{\text{out}} + S^\lambda_{\text{in}} - C^r_{\text{out}} + C^r_{\text{in}}
\end{equation}

各项分别表示：
\begin{itemize}
    \item $S^\lambda_{\text{out}}$：由于流动导致的概率流出
    \item $S^\lambda_{\text{in}}$：由于流动导致的概率流入
    \item $C^r_{\text{out}}$：由于转换导致的概率流出
    \item $C^r_{\text{in}}$：由于转换导致的概率流入
\end{itemize}

\subsubsection{概率流的具体表达式}

\textbf{流动项}（出流）：
\begin{align}
S^\lambda_{\text{out}} &= \sum_i \left[\lambda n_i^+ + \lambda n_i^-\right] P \label{eq:stream-out}\\
C^r_{\text{out}} &= \sum_i \left[r_+ n_i^+ + r_- n_i^-\right] P \label{eq:conv-out}
\end{align}

\textbf{流动项}（入流，来自左右邻居的贡献）：
\begin{align}
S^\lambda_{\text{in}} &= \sum_i \Big[\lambda(n_{i-1}^+ + 1)P(\ldots, n_{i-1}^+ + 1, n_i^+ - 1, \ldots) \nonumber\\
&\qquad + \lambda(n_{i+1}^- + 1)P(\ldots, n_i^- - 1, n_{i+1}^- + 1, \ldots)\Big] \label{eq:stream-in}
\end{align}

\textbf{转换项}（入流）：
\begin{align}
C^r_{\text{in}} &= \sum_i \Big[r_+(n_i^+ + 1, n_i^- - 1)(n_i^+ + 1)P(\ldots, n_i^+ + 1, n_i^- - 1, \ldots) \nonumber\\
&\qquad + r_-(n_i^+ - 1, n_i^- + 1)(n_i^- + 1)P(\ldots, n_i^+ - 1, n_i^- + 1, \ldots)\Big] \label{eq:conv-in}
\end{align}

\subsection{Poisson 表示法}\label{subsec:poisson-rep}

直接求解 Master 方程是不可行的，因为对于相互作用系统（即非平凡的转换率），右边是非线性的。因此，我们采用\textbf{Poisson 表示法}\cite{gardiner1977}，这是一种将离散的随机过程转化为连续过程的方法。

\subsubsection{核心思想}

\begin{notebox}[Poisson 表示法的核心思想]
将离散的占据数 $n_i^\pm$（正整数上的随机过程）替换为连续的 Poisson 速率 $\alpha_i, \beta_i$（分别对应 $+$ 和 $-$ 粒子）。这些速率作为 Poisson 过程的参数，然后生成占据数 $n_i^\pm$。

统计矩通过对 Poisson 过程在固定场下的实现以及这些场的分布进行平均来获得。
\end{notebox}

概率分布 $P(\{n_i^+, n_i^-\})$ 与伪分布 $f(\{\alpha_i, \beta_i\})$ 的关系为：
\begin{equation}\label{eq:poisson-rep}
\boxed{P = \int \prod_i \dd\alpha_i \, \dd\beta_i \; e^{-\alpha_i}\frac{\alpha_i^{n_i^+}}{n_i^+!} \cdot e^{-\beta_i}\frac{\beta_i^{n_i^-}}{n_i^-!} \; f}
\end{equation}

\subsubsection{从一般离散过程到 Poisson 过程的不匹配}

由于一般离散过程的通用性与局部 Poisson 过程之间存在``不匹配''，需要更高层次的抽象：
\begin{enumerate}
    \item 局部速率 $\alpha_i, \beta_i$ 是\textbf{复数}
    \item 分布 $f$ 实际上是一个\textbf{伪分布}，因为它也可以是局部负的
\end{enumerate}

这些扩展并非非物理的，因为 Poisson 场是不可直接观测的辅助数学对象。\textbf{可观测}的是物理过程的矩，它们是正且实的。

\begin{remark}[虚数的直观理解]
虚数对 Poisson 场的贡献可以看作是偏离 Poisson 统计的表现，这种偏离显式地耦合了随机变量的均值和方差。

比较一个速率为 $\lambda$ 的 Poisson 变量 $n$ 和一个速率为 $\lambda + i\xi$ 的变量 $m$，其中 $\xi$ 是均值为零的随机变量，$\langle \xi^2 \rangle = c$。对于 Poisson 过程，有：
\[
\langle n \rangle = \langle n^2 \rangle_c = \lambda
\]
但对于 $m$：
\[
\langle m \rangle = \lambda, \quad \langle m^2 \rangle_c = \lambda - c^2
\]
虚噪声导致微观分布的\textbf{变窄}(narrowing)，而实噪声会导致\textbf{变宽}(widening)。
\end{remark}

\subsubsection{替换规则：从 Master 方程到 Fokker-Planck 方程}

通过将 Poisson 表示代入 Master 方程，并利用分部积分，我们得到以下替换规则。以下省略 $f$ 和 $P$ 中与前述默认参数一致的参数。

\begin{derivation}[线性项的替换规则（流动和自发转换）]
\textbf{出流项}（右移）：
\begin{equation}\label{eq:replace-out-right}
n_i^+ P \;\to\; (\alpha_i f - \partial_{\alpha_i}\alpha_i f)
\end{equation}

\textbf{出流项}（左移）：
\begin{equation}\label{eq:replace-out-left}
n_i^- P \;\to\; (\beta_i f - \partial_{\beta_i}\beta_i f)
\end{equation}

\textbf{入流项}（从左边流入）：
\begin{equation}\label{eq:replace-in-left}
(n_{i-1}^+ + 1)P(\ldots, n_{i-1}^+ + 1, n_i^+ - 1, \ldots) \;\to\; (\alpha_{i-1} - \partial_{\alpha_i}\alpha_{i-1})f
\end{equation}

\textbf{入流项}（从右边流入）：
\begin{equation}\label{eq:replace-in-right}
(n_{i+1}^- + 1)P(\ldots, n_i^- - 1, n_{i+1}^- + 1, \ldots) \;\to\; (\beta_{i+1} - \partial_{\beta_i}\beta_{i+1})f
\end{equation}
\end{derivation}

\begin{derivation}[非线性项的替换规则（转换过程）]
\textbf{$(+)$ 到 $(-)$ 的转换}：
\begin{equation}\label{eq:replace-conv-plus}
n_i^+(n_i^+ - 1)P \;\to\; (1 - \partial_{\alpha_i})^2 \alpha_i^2 f
\end{equation}

\textbf{$(-)$ 到 $(+)$ 的转换}：
\begin{equation}\label{eq:replace-conv-minus}
n_i^-(n_i^- - 1)P \;\to\; (1 - \partial_{\beta_i})^2 \beta_i^2 f
\end{equation}

\textbf{$(+)$ 粒子因 $(-)$ 粒子而转换}：
\begin{equation}\label{eq:replace-conv-mix-plus}
(n_i^- + 1)P(\ldots, n_i^+ - 1, n_i^- + 1, \ldots) \;\to\; (-1 + \partial_{\alpha_i})(-1 + \partial_{\beta_i})\alpha_i^2 f
\end{equation}

\textbf{$(-)$ 粒子因 $(+)$ 粒子而转换}：
\begin{equation}\label{eq:replace-conv-mix-minus}
(n_i^+ + 1)P(\ldots, n_i^+ + 1, n_i^- - 1, \ldots) \;\to\; (-1 + \partial_{\alpha_i})(-1 + \partial_{\beta_i})\beta_i^2 f
\end{equation}
\end{derivation}

\subsubsection{详细推导：替换规则的获得}

以式~(\ref{eq:replace-out-right})为例，展示详细推导过程：
\begin{align}
n_i^+ P &= \int \prod_j \dd\alpha_j \, \dd\beta_j \; n_i^+ \frac{e^{-\alpha_i}\alpha_i^{n_i^+}}{n_i^+!} \cdot \frac{e^{-\beta_i}\beta_i^{n_i^-}}{n_i^-!} \; f \nonumber\\
&= \int \prod_j \dd\alpha_j \, \dd\beta_j \; \frac{e^{-\alpha_i}\alpha_i^{n_i^+}}{(n_i^+ - 1)!} \cdot \frac{e^{-\beta_i}\beta_i^{n_i^-}}{n_i^-!} \; f \nonumber\\
&= \int \prod_j \dd\alpha_j \, \dd\beta_j \; \frac{e^{-\alpha_i}\alpha_i^{n_i^+}}{n_i^+!} \cdot \frac{e^{-\beta_i}\beta_i^{n_i^-}}{n_i^-!} \; \alpha_i f \nonumber\\
&\quad - \int \prod_j \dd\alpha_j \, \dd\beta_j \; (-\partial_{\alpha_i})\left[e^{-\alpha_i}\frac{\alpha_i^{n_i^+}}{n_i^+!}\right] \cdot \frac{e^{-\beta_i}\beta_i^{n_i^-}}{n_i^-!} \; \alpha_i f \nonumber\\
&= P \text{ 中 } n_i^+ \text{ 被替换后的形式} \nonumber
\end{align}

经过分部积分（假设边界项消失），得到替换规则。

对于非线性项，例如式~(\ref{eq:replace-conv-plus})：
\begin{align}
n_i^+(n_i^+ - 1)P &= \int \dd\alpha_i \; n_i^+(n_i^+ - 1) \frac{e^{-\alpha_i}\alpha_i^{n_i^+}}{n_i^+!} \; f \nonumber\\
&= \int \dd\alpha_i \; \frac{e^{-\alpha_i}\alpha_i^{n_i^+}}{(n_i^+ - 2)!} \; f \nonumber\\
&= \int \dd\alpha_i \; (1 - \partial_{\alpha_i})^2 \left[\frac{e^{-\alpha_i}\alpha_i^{n_i^+}}{n_i^+!}\right] \alpha_i^2 f \nonumber\\
&\to\; (1 - \partial_{\alpha_i})^2 \alpha_i^2 f \nonumber
\end{align}

\subsection{Fokker-Planck 方程}\label{subsec:fokker-planck}

将上述替换规则收集起来，得到关于伪概率密度 $f(\{\alpha_i, \beta_i\})$ 的 Fokker-Planck 方程：
\begin{equation}\label{eq:fpe}
\boxed{\partial_t f = F_\lambda + F_r + J}
\end{equation}

其中各项分别为：

\subsubsection{流动贡献 $F_\lambda$}
\begin{align}
F_\lambda &= \lambda \sum_i \Big[(\alpha_{i-1} - \alpha_i) + \partial_{\alpha_i}(\alpha_i - \alpha_{i-1})\Big]f \nonumber\\
&\quad - \lambda \sum_i \Big[(\beta_i - \beta_{i+1}) + \partial_{\beta_i}(\beta_{i+1} - \beta_i)\Big]f \label{eq:f-lambda}
\end{align}

\subsubsection{转换贡献 $F_r$}
\begin{align}
F_r &= r_0 \sum_i \Big[\partial_{\alpha_i}(\alpha_i - \beta_i) - \partial_{\beta_i}(\alpha_i - \beta_i)\Big]f \nonumber\\
&\quad + r_1 \sum_i \Big[\partial_{\alpha_i} - \partial_{\beta_i}\Big](\alpha_i + \beta_i)(\alpha_i - \beta_i)f \label{eq:f-r}
\end{align}

\subsubsection{噪声项 $J$}

\begin{equation}\label{eq:noise-term}
\boxed{J = -r_1 \sum_i \Big[\partial_{\alpha_i}^2 \alpha_i^2 + \partial_{\beta_i}^2 \beta_i^2\Big]f}
\end{equation}

最后一行包含二阶导数项，这在对应的 Langevin 方程中表现为\textbf{噪声}。注意这个噪声项的系数是\textbf{负}的，这最终导致了虚噪声的出现。

\subsection{确定性动力学方程}\label{subsec:det-eom}

从 Fokker-Planck 方程中，我们可以读出确定性方程（忽略噪声项）。将记号从 $\lambda$ 改为 $v$（与传统一致），得到：
\begin{align}
\partial_t \alpha\big|_{\text{det}} &= -v\nabla\alpha - r_0(\alpha - \beta) - r_1(\alpha + \beta)(\alpha - \beta) \label{eq:det-alpha}\\
\partial_t \beta\big|_{\text{det}} &= v\nabla\beta + r_0(\alpha - \beta) + r_1(\alpha + \beta)(\alpha - \beta) \label{eq:det-beta}
\end{align}

注意这里的 $\nabla$ 是最近邻差分的简写，\textbf{尚未}取连续极限。

\subsection{噪声分析}\label{subsec:noise-analysis}

为了确定噪声项，我们将 Fokker-Planck 方程写成标准形式：
\begin{equation}\label{eq:fpe-standard}
\partial_t f = -\sum_i \partial_{x_i}(\mu_i f) + \sum_{i,j} \partial_{x_i}\partial_{x_j}[D_{ij} f]
\end{equation}

扩散矩阵 $D_{ij}$ 为：
\begin{equation}\label{eq:diffusion-matrix}
D_{ij} = -r_1 \begin{pmatrix} \alpha_i^2 & 0 \\ 0 & \beta_i^2 \end{pmatrix}
\end{equation}

引入变量替换：
\begin{equation}\label{eq:variable-change}
\tilde{\rho}_i = \alpha_i + \beta_i \quad (\text{总密度}), \qquad \tilde{m}_i = \alpha_i - \beta_i \quad (\text{磁化强度})
\end{equation}

扩散矩阵变为：
\begin{equation}\label{eq:diffusion-rho-m}
D_{ij} = -r_1 \begin{pmatrix} \frac{\tilde{\rho}^2 + \tilde{m}^2}{4} & \frac{-2\tilde{\rho}\tilde{m}}{4} \\ \frac{-2\tilde{\rho}\tilde{m}}{4} & \frac{\tilde{\rho}^2 + \tilde{m}^2}{4} \end{pmatrix}
\end{equation}

\subsubsection{矩阵平方根与虚噪声}

将 Fokker-Planck 方程翻译为 Langevin 方程需要计算矩阵平方根 $B = \sqrt{2D}$。在物理相关的极限 $\tilde{\rho} \approx \rho_0$, $\tilde{m} \approx 0$ 下，到 $\tilde{m}$ 的零阶：
\begin{equation}\label{eq:matrix-sqrt}
B = i\rho_0\sqrt{\frac{r_1}{2}} \begin{pmatrix} 1 & -1 \\ -1 & 1 \end{pmatrix}
\end{equation}

这里出现了因子 $i = \sqrt{-1}$！这是 Poisson 表示法的标志性结果。

\begin{notebox}[虚噪声的物理意义]
两个特征值为：
\[
\lambda_\rho = 0, \qquad \lambda_m = i\sqrt{2r_1\rho_0}
\]
$\lambda_\rho = 0$ 意味着\textbf{密度场没有噪声}，而磁化强度场有虚噪声。

这并不意味着任何可观测量是虚的。虚噪声的出现标志着偏离 Poisson 统计——\textbf{在 Poisson 表示法中，只有复数场才能描述非 Poisson 统计}，因为 Poisson 分布的均值和方差必须相等，而要打破这一等式，场的虚部必须非零。
\end{notebox}

\subsection{Langevin 方程}\label{subsec:langevin}

总结以上分析，密度场和磁化强度场的耦合 Langevin 方程为：

\begin{eqbox}[Langevin 方程组]
\begin{align}
\partial_t \tilde{\rho} &= -v\nabla\tilde{m} \label{eq:langevin-rho}\\
\partial_t \tilde{m} &= -v\nabla\tilde{\rho} - r_0\tilde{m} - r_1\tilde{\rho}\tilde{m} + i\sqrt{2r_1\tilde{\rho}}\,\xi \label{eq:langevin-m}
\end{align}
\end{eqbox}

其中 $\xi$ 是高斯不相关噪声，满足：
\begin{equation}\label{eq:noise-correlation}
\langle \xi \rangle = 0, \qquad \langle \xi(t)\xi(t') \rangle = \delta(t - t')
\end{equation}

\begin{remark}[与已有文献的联系]
这组方程的确定性版本（设 $r_0 \equiv 0$）与反排列粒子的闭合流体动力学\cite{boltz2025}在一维设定下的结构相同。一般的 Toner-Tu 型方法\cite{toner1998}也涵盖这些动力学，但不对 $r_0$ 和 $r_1$ 做出具体预测。
\end{remark}

\subsection{结构因子的解析计算}\label{subsec:structure-factor}

\subsubsection{线性化与简化}

考虑围绕均匀非极化态的小扰动：
\begin{equation}\label{eq:linearization}
\tilde{\rho} = \rho_0 + \delta\tilde{\rho}, \qquad \tilde{m} \text{ 作为小量处理}
\end{equation}

将乘性噪声近似为加性噪声（忽略 $\delta\tilde{\rho}\,\xi$ 项），得到闭的二阶方程：
\begin{equation}\label{eq:second-order-eq}
\partial_t^2 \delta\tilde{\rho} = v^2\nabla^2\delta\tilde{\rho} - (r_0 + r_1\rho_0)\partial_t\delta\tilde{\rho} - i\sqrt{2r_1\rho_0}\,v\nabla\xi
\end{equation}

\subsubsection{过阻尼近似}

在 Fourier 空间（波数为 $k$）中，确定性部分构成一个谐振子。在 $k \to 0$ 的极限下，它是\textbf{过阻尼的}：两个动力学特征值都是负的。长时间行为由较大的特征值主导：
\begin{equation}\label{eq:eigenvalue}
\mu = -\frac{v^2 k^2}{r_0 + r_1\rho_0}
\end{equation}

这给出了扩散标度 $\sim e^{-Dk^2t}$，其中扩散系数：
\begin{equation}\label{eq:diffusivity}
D = \frac{v^2}{r_0 + r_1\rho_0}
\end{equation}

\subsubsection{一阶 Langevin 方程}

由于确定性动力学在大尺度上是过阻尼的，可以忽略惯性项，将所有内容约化为一个一阶 Langevin 方程（或噪声扩散方程）：
\begin{equation}\label{eq:first-order-langevin}
\partial_t \delta\tilde{\rho}_k = -Dk^2\delta\tilde{\rho}_k + ic\eta_k
\end{equation}

其中：
\begin{equation}\label{eq:noise-strength}
c = v\sqrt{\frac{2r_1\rho_0}{r_0 + r_1\rho_0}}
\end{equation}

\subsubsection{关联噪声的定义}

为了正确理解 Fourier 空间中的有效噪声统计，需要注意 $\nabla\xi$ 实际上是 $\xi_{i+1} - \xi_i$。为此，引入无漂移的关联噪声 $\eta_i = \eta(i)$：
\begin{equation}\label{eq:correlated-noise}
\langle \eta_i(t)\eta_j(t') \rangle = \delta(t - t') \times \begin{cases} 2, & i = j \\ -1, & i = j \pm 1 \\ 0, & \text{其他} \end{cases}
\end{equation}

在 Fourier 空间：
\begin{equation}\label{eq:noise-fourier}
\langle \eta_k(t)\eta_{k'}(t') \rangle = \delta(k + k')\delta(t - t') \cdot 2[1 - \cos(k)]
\end{equation}

\subsubsection{结构因子的最终表达式}

利用上述关联噪声求解方程~(\ref{eq:first-order-langevin})，得到大时间极限下的结构因子。注意结构因子可以通过关联 Poisson 场的涨落来计算\cite{bertrand2019}：
\begin{equation}\label{eq:ns-definition}
N_S(k) = \langle \rho_k \rho_{-k} \rangle = \rho_0 + \langle \delta\tilde{\rho}_k \, \delta\tilde{\rho}_{-k} \rangle
\end{equation}

\begin{derivation}[结构因子的详细计算]
方程~(\ref{eq:first-order-langevin})的解可以通过 Fourier 变换获得。设 $\delta\tilde{\rho}_k(t)$ 的稳态涨落为：
\[
\delta\tilde{\rho}_k(t) = ic \int_{-\infty}^{t} e^{-Dk^2(t-t')}\eta_k(t')\,\dd t'
\]

因此：
\begin{align}
\langle \delta\tilde{\rho}_k \, \delta\tilde{\rho}_{-k} \rangle &= -c^2 \int_{-\infty}^{t}\int_{-\infty}^{t} e^{-Dk^2(2t - t' - t'')} \langle \eta_k(t')\eta_{-k}(t'') \rangle \, \dd t' \, \dd t'' \nonumber\\
&= -c^2 \int_{-\infty}^{t} e^{-2Dk^2(t-t')} \cdot 2[1 - \cos(k)] \, \dd t' \nonumber\\
&= -c^2 \frac{2[1 - \cos(k)]}{2Dk^2} \nonumber\\
&= -\frac{c^2[1 - \cos(k)]}{Dk^2}\nonumber
\end{align}
\end{derivation}

最终结果为：
\begin{resultbox}[结构因子（核心结果）]
\begin{equation}\label{eq:structure-factor-final}
\boxed{N_S(k) = \rho_0 - \frac{c^2[1 - \cos(k)]}{Dk^2} = \rho_0 - \frac{c^2}{2D}\left(1 - \frac{k^2}{12} + \ldots\right)}
\end{equation}
\end{resultbox}

\subsubsection{物理分析}

比较 $S_0 = \lim_{k\to 0^+} N_S(k) = \rho_0 - c^2/(2D)$ 与 Poisson 值 $S|_{r_1=0} = \rho_0$：
\begin{equation}\label{eq:reduction-ratio}
\frac{S_0}{\rho_0} = 1 - \frac{c^2}{2D\rho_0} = 1 - \frac{r_1\rho_0}{r_0 + r_1\rho_0}
\end{equation}

在相互作用主导的极限 $r_1\rho_0 \gg r_0$ 下：
\begin{equation}\label{eq:limit-result}
\boxed{\frac{S_0}{\rho_0} \to \frac{1}{2}}
\end{equation}

这意味着长程密度涨落被抑制到 Poisson 值的一半！

\begin{notebox}[关键物理结论]
\begin{enumerate}
    \item 当 $r_1\rho_0 \gg r_0$（强反排列相互作用）时，$S_0/\rho_0 \to 1/2$
    \item 引入自发转换 $r_0$ 等效于加入取向噪声，支持了``动态建立的关联在确定性模型中更为显著''的早期论断
    \item 更低值（包括超均匀的 $S_0 \to 0$）可能在质心或极化严格守恒的模型中实现
    \item 在 $k \to 0$ 的极限下，结构因子趋近于一个\textbf{有限常数}而非零，因此系统\textbf{不是严格超均匀的}，但展现了显著的涨落抑制
\end{enumerate}
\end{notebox}

\subsection{数值验证（一维模型）}\label{subsec:1d-numerics}

论文还报告了一维模型的数值模拟结果。使用 Gillespie 算法\cite{boltz2025}直接模拟随机过程，模拟参数为 $\lambda = 1$, $r_1 = 1/100$，系统大小 $L$ 可变。

\begin{itemize}
    \item 模拟结果与方程~(\ref{eq:structure-factor-final})的预测在定性上一致
    \item 对于足够大的系统尺寸，显示出长程密度涨落的抑制
    \item 需要注意在有限系统中，Binomial 统计与 Poisson 统计之间存在差异，需要进行修正
\end{itemize}

\begin{remark}[有限尺寸效应与虚假超均匀性]
如果在小 $k$ 范围内忽略或不了解偏移量 $S_0$，可能会导致错误的推论——即推断出一个虚假的幂律行为。在图~2 的情况下，这可能导致推断 $k^{\approx 0.25}$ 的``反常''标度。这提醒我们：在没有主要理论基础的情况下，从数值数据推断超均匀性是有风险的。
\end{remark}
